\documentclass{article}
\usepackage{amsmath,amssymb,amsthm}
\usepackage{algorithm}
\usepackage{algpseudocode}
\usepackage{hyperref}
\usepackage{geometry}
\geometry{margin=1in}
\newtheorem{theorem}{Theorem}
\newtheorem{lemma}{Lemma}
\newtheorem{assumption}{Assumption}
\newcommand{\E}{\mathbb{E}}
\newcommand{\norm}[1]{\left\|#1\right\|}
\newcommand{\inner}[1]{\left\langle#1\right\rangle}
\begin{document}

\section*{Algorithm Name}
\textbf{Federated Averaging with Local SGD (FedAvg‑LSGD).}  
A set of $K$ clients collaboratively minimize a global non‑convex smooth objective  

\[
f(w)\;:=\;\frac{1}{K}\sum_{k=1}^{K}F_k(w),\qquad 
F_k(w)\;=\;\mathbb{E}_{\xi\sim\mathcal{D}_k}\big[\ell(w;\xi)\big],
\]

where $F_k$ is the local empirical risk on client~$k$.

\section*{Mathematical Formulation}
\begin{itemize}
    \item \textbf{Hyper‑parameters}
    \begin{itemize}
        \item Global learning rate $\eta>0$ (used locally as well).
        \item Number of local SGD steps per communication round $E\in\mathbb{N}$.
        \item Total number of communication rounds $T$.
    \end{itemize}
    \item \textbf{Initialization}
    \[
    w^{0}\in\mathbb{R}^{d}\quad\text{(central model).}
    \]
    \item \textbf{Local update on client $k$ (for $e=0,\dots,E-1$)}
    \[
    w^{t,e+1}_{k}=w^{t,e}_{k}-\eta\, g^{t,e}_{k},
    \qquad 
    g^{t,e}_{k}= \nabla \ell\big(w^{t,e}_{k};\xi^{t,e}_{k}\big),
    \]
    where $\xi^{t,e}_{k}\sim\mathcal{D}_k$ is drawn independently.
    \item \textbf{Model averaging (communication) after $E$ local steps}
    \[
    w^{t+1}\;=\;\frac{1}{K}\sum_{k=1}^{K} w^{t,E}_{k}.
    \]
    \item \textbf{Drift term}
    Define the \emph{client drift} at round $t$
    \[
    \Delta^{t}_{k}\;:=\;w^{t,E}_{k}-w^{t},
    \]
    which captures the deviation of a client’s local model from the global model after $E$ steps.
\end{itemize}

\section*{Pseudocode}
\begin{algorithm}[H]
\caption{FedAvg‑LSGD}
\begin{algorithmic}[1]
\State \textbf{Input:} learning rate $\eta$, local steps $E$, rounds $T$, initial model $w^{0}$.
\For{$t=0$ \textbf{to} $T-1$}
    \State Server broadcasts $w^{t}$ to all $K$ clients.
    \For{each client $k\in\{1,\dots,K\}$ \textbf{in parallel}}
        \State $w^{t,0}_{k}\gets w^{t}$
        \For{$e=0$ \textbf{to} $E-1$}
            \State Sample $\xi^{t,e}_{k}\sim\mathcal{D}_{k}$
            \State $g^{t,e}_{k}\gets \nabla\ell\big(w^{t,e}_{k};\xi^{t,e}_{k}\big)$
            \State $w^{t,e+1}_{k}\gets w^{t,e}_{k}-\eta\, g^{t,e}_{k}$
        \EndFor
        \State Send $w^{t,E}_{k}$ to the server.
    \EndFor
    \State Server aggregates: $w^{t+1}\gets \frac{1}{K}\sum_{k=1}^{K} w^{t,E}_{k}$
\EndFor
\State \textbf{Output:} $w^{T}$.
\end{algorithmic}
\end{algorithm}

\section*{Dry Run Example}
Consider the scalar quadratic $f(w)=(w-3)^{2}$, thus $\nabla f(w)=2(w-3)$.  
Take a single client ($K=1$), learning rate $\eta=0.1$, $E=1$ (one local step per round), and start from $w^{0}=0$.

\begin{center}
\begin{tabular}{c|c|c|c}
Round $t$ & $w^{t}$ & $g^{t,0}= \nabla\ell(w^{t};\cdot)$ & $w^{t+1}=w^{t}-\eta g^{t,0}$\\ \hline
0 & $0$ & $2(0-3)=-6$ & $0-0.1(-6)=0.6$\\
1 & $0.6$ & $2(0.6-3)=-4.8$ & $0.6-0.1(-4.8)=1.08$\\
2 & $1.08$ & $2(1.08-3)=-3.84$ & $1.08-0.1(-3.84)=1.464$\\
\end{tabular}
\end{center}
Objective values: $f(0)=9$, $f(0.6)=5.76$, $f(1.08)=3.6864$, $f(1.464)=2.3545$.  
The algorithm reduces the loss despite the non‑convex (here actually convex) landscape.

\newpage
\section*{Assumptions}
\begin{assumption}[Smoothness]
Each local function $F_k$ is $L$‑smooth:
\[
\forall w,v\in\mathbb{R}^{d},\quad
F_k(v)\le F_k(w)+\inner{\nabla F_k(w)}{v-w}+\frac{L}{2}\norm{v-w}^{2}.
\] 
\end{assumption}
\begin{assumption}[Unbiased Stochastic Gradients and Bounded Variance]
For any client $k$, round $t$ and local step $e$,
\[
\E_{\xi^{t,e}_{k}}\!\big[g^{t,e}_{k}\mid w^{t,e}_{k}\big]=\nabla F_k\big(w^{t,e}_{k}\big),\qquad
\E_{\xi^{t,e}_{k}}\!\big[\|g^{t,e}_{k}-\nabla F_k(w^{t,e}_{k})\|^{2}\mid w^{t,e}_{k}\big]\le\sigma^{2}.
\]
\end{assumption}
\begin{assumption}[Bounded Gradient Norm]
\[
\forall k,\; \forall w,\quad \|\nabla F_k(w)\|\le G.
\]
\end{assumption}
\begin{assumption}[Client Dissimilarity (Drift) \cite{li2020federated}]
There exists $\delta\ge0$ such that for all $w$,
\[
\frac{1}{K}\sum_{k=1}^{K}\|\nabla F_k(w)-\nabla f(w)\|^{2}\le\delta^{2}.
\]
\end{assumption}
All constants $L,G,\sigma,\delta$ are finite and independent of $t$.

\section*{Main Convergence Theorem}
\begin{theorem}[Non‑convex convergence of FedAvg‑LSGD]\label{thm:main}
Under Assumptions 1–4, let the learning rate satisfy $0<\eta\le\frac{1}{2L(E+1)}$. After $T$ communication rounds (i.e. $TE$ local updates) the iterates produced by FedAvg‑LSGD obey  

\[
\frac{1}{T}\sum_{t=0}^{T-1}\E\!\big[\|\nabla f(w^{t})\|^{2}\big]
\;\le\;
\underbrace{\frac{2\big(f(w^{0})-f^{\star}\big)}{\eta TE}}_{\text{optimization term}}
\;+\;
\underbrace{4L^{2}\eta^{2}E^{2}G^{2}}_{\text{drift term}}
\;+\;
\underbrace{\frac{2\sigma^{2}}{K\eta TE}}_{\text{stochastic variance}} .
\]
Consequently, choosing $\eta = \Theta\!\big(\frac{1}{\sqrt{TE}}\big)$ yields  

\[
\frac{1}{T}\sum_{t=0}^{T-1}\E\!\big[\|\nabla f(w^{t})\|^{2}\big]
= O\!\Big(\frac{1}{\sqrt{TE}}+L^{2}E^{2}\frac{1}{TE}\Big)
= O\!\Big(\frac{1}{\sqrt{TE}}+L^{2}E^{2}\frac{1}{TE}\Big).
\]
\end{theorem}

\section*{Theoretical Properties}
\begin{itemize}
    \item \textbf{Stability w.r.t. $\eta$.} The bound requires $\eta\le 1/(2L(E+1))$; larger $\eta$ would break the smoothness descent inequality.
    \item \textbf{Effect of local steps $E$.} The drift term scales as $L^{2}\eta^{2}E^{2}G^{2}$, showing a quadratic penalty for too many local updates.
    \item \textbf{Client heterogeneity.} When $\delta=0$ (identical clients) the drift term can be refined to $L^{2}\eta^{2}E^{2}\delta^{2}$, which vanishes, recovering the classical $O(1/\sqrt{TE})$ rate of centralized SGD.
    \item \textbf{Comparison.} For $E=1$ the theorem reduces to the standard non‑convex SGD rate $O(1/\sqrt{T})$ (up to the $1/K$ variance reduction). For $E>1$ the extra drift term explains the empirical trade‑off between communication efficiency and convergence speed.
\end{itemize}

\section*{Discussion}
The theorem demonstrates that FedAvg‑LSGD attains the same $O(1/\sqrt{TE})$ convergence as centralized SGD provided the number of local steps $E$ grows slower than $\sqrt{T}$, i.e. $E=O(T^{1/4})$. This matches known practical guidelines: moderate local computation reduces communication without sacrificing statistical efficiency.

\newpage
\section*{Preliminaries (Definitions and Lemmas)}
\begin{lemma}[Smoothness descent]\label{lem:smooth}
If $F_k$ is $L$‑smooth, then for any $u,v\in\mathbb{R}^{d}$,
\[
F_k(v)\le F_k(u)+\inner{\nabla F_k(u)}{v-u}+\frac{L}{2}\norm{v-u}^{2}.
\] 
\end{lemma}
\begin{lemma}[One‑step SGD bound]\label{lem:sgd}
For a single SGD step $w^{+}=w-\eta g$ with $g$ unbiased and $\E\|g-\nabla F_k(w)\|^{2}\le\sigma^{2}$,
\[
\E\big[F_k(w^{+})\big]\le F_k(w)-\eta\|\nabla F_k(w)\|^{2}
+\frac{L\eta^{2}}{2}\big(\|\nabla F_k(w)\|^{2}+\sigma^{2}\big).
\] 
\end{lemma}
\begin{proof}
Apply Lemma~\ref{lem:smooth} with $u=w$, $v=w^{+}=w-\eta g$,
take expectation conditioned on $w$, use unbiasedness and variance bound, and simplify. \hfill $\square$
\end{proof}

\section*{Main Proof (Step‑by‑Step Derivation)}
We prove Theorem~\ref{thm:main}. Let $w^{t}$ be the global model at the beginning of round $t$.
For each client $k$ define the local iterates $w^{t,0}_{k}=w^{t}$ and for $e=0,\dots,E-1$  

\[
w^{t,e+1}_{k}=w^{t,e}_{k}-\eta g^{t,e}_{k},\qquad 
g^{t,e}_{k}= \nabla\ell(w^{t,e}_{k};\xi^{t,e}_{k}).
\tag{1}
\]

The global model after aggregation is  

\[
w^{t+1}= \frac{1}{K}\sum_{k=1}^{K} w^{t,E}_{k}.
\tag{2}
\]

\subsection*{Step 1: Descent of a single local step}
From Lemma~\ref{lem:sgd} applied to client $k$ at local step $e$,
\[
\E\big[F_k(w^{t,e+1}_{k})\mid w^{t,e}_{k}\big]
\le
F_k(w^{t,e}_{k})-\eta\|\nabla F_k(w^{t,e}_{k})\|^{2}
+\frac{L\eta^{2}}{2}\big(\|\nabla F_k(w^{t,e}_{k})\|^{2}+\sigma^{2}\big).
\tag{3}
\]

\subsection*{Step 2: Unrolling $E$ local steps}
Summing (3) over $e=0,\dots,E-1$ and taking full expectation yields

\[
\begin{aligned}
\E\big[F_k(w^{t,E}_{k})\big]
&\le
F_k(w^{t})-\eta\sum_{e=0}^{E-1}\E\!\big[\|\nabla F_k(w^{t,e}_{k})\|^{2}\big]
+\frac{L\eta^{2}}{2}\sum_{e=0}^{E-1}\E\!\big[\|\nabla F_k(w^{t,e}_{k})\|^{2}+\sigma^{2}\big]  \\
&=F_k(w^{t})-(\eta-\tfrac{L\eta^{2}}{2})\sum_{e=0}^{E-1}\E\!\big[\|\nabla F_k(w^{t,e}_{k})\|^{2}\big]
+\frac{L\eta^{2}E\sigma^{2}}{2}.
\end{aligned}
\tag{4}
\]

\subsection*{Step 3: Relating local gradients to the global gradient}
Add and subtract $\nabla f(w^{t})$ inside the norm and use the inequality 
$\|a+b\|^{2}\le2\|a\|^{2}+2\|b\|^{2}$:

\[
\|\nabla F_k(w^{t,e}_{k})\|^{2}
\le 2\|\nabla F_k(w^{t})\|^{2}+2\|\nabla F_k(w^{t,e}_{k})-\nabla F_k(w^{t})\|^{2}.
\tag{5}
\]

By $L$‑smoothness,
\[
\|\nabla F_k(w^{t,e}_{k})-\nabla F_k(w^{t})\|
\le L\norm{w^{t,e}_{k}-w^{t}}.
\tag{6}
\]

Using the definition of drift $\Delta^{t}_{k}=w^{t,E}_{k}-w^{t}$ and the recursion (1), we obtain

\[
w^{t,e}_{k}-w^{t}= -\eta\sum_{s=0}^{e-1} g^{t,s}_{k},
\qquad
\|\!w^{t,e}_{k}-w^{t}\!\|\le \eta\sum_{s=0}^{e-1}\|g^{t,s}_{k}\|.
\tag{7}
\]

With the bounded gradient norm $\|g^{t,s}_{k}\|\le G$ (Assumption 3) we have  

\[
\|\!w^{t,e}_{k}-w^{t}\!\|\le \eta e G\le \eta EG .
\tag{8}
\]

Combining (6)–(8),

\[
\|\nabla F_k(w^{t,e}_{k})-\nabla F_k(w^{t})\|^{2}\le L^{2}\eta^{2}E^{2}G^{2}.
\tag{9}
\]

Insert (9) into (5) and then into (4):

\[
\begin{aligned}
\E\big[F_k(w^{t,E}_{k})\big]
&\le
F_k(w^{t})-(\eta-\tfrac{L\eta^{2}}{2})\sum_{e=0}^{E-1}\Big(2\E\!\big[\|\nabla F_k(w^{t})\|^{2}\big]
-2L^{2}\eta^{2}E^{2}G^{2}\Big)\\
&\quad+\frac{L\eta^{2}E\sigma^{2}}{2}.
\end{aligned}
\tag{10}
\]

Since the sum over $e$ contributes a factor $E$,

\[
\E\big[F_k(w^{t,E}_{k})\big]\le
F_k(w^{t})-2E(\eta-\tfrac{L\eta^{2}}{2})\E\!\big[\|\nabla F_k(w^{t})\|^{2}\big]
+2E(\eta-\tfrac{L\eta^{2}}{2})L^{2}\eta^{2}E^{2}G^{2}
+\frac{L\eta^{2}E\sigma^{2}}{2}.
\tag{11}
\]

\subsection*{Step 4: Averaging over clients}
Average (11) over $k=1,\dots,K$ and use the definition of the global function $f$:

\[
\begin{aligned}
\E\big[f(w^{t,E})\big]
&:=\frac{1}{K}\sum_{k=1}^{K}\E\big[F_k(w^{t,E}_{k})\big] \\
&\le f(w^{t})-2E(\eta-\tfrac{L\eta^{2}}{2})\underbrace{\frac{1}{K}\sum_{k=1}^{K}\E\!\big[\|\nabla F_k(w^{t})\|^{2}\big]}_{\text{decompose via drift}}\\
&\quad+2E(\eta-\tfrac{L\eta^{2}}{2})L^{2}\eta^{2}E^{2}G^{2}
+\frac{L\eta^{2}E\sigma^{2}}{2}.
\end{aligned}
\tag{12}
\]

Decompose the first term using Assumption~4:

\[
\frac{1}{K}\sum_{k=1}^{K}\|\nabla F_k(w^{t})\|^{2}
= \|\nabla f(w^{t})\|^{2}
+ \frac{1}{K}\sum_{k=1}^{K}\|\nabla F_k(w^{t})-\nabla f(w^{t})\|^{2}
\le \|\nabla f(w^{t})\|^{2}+\delta^{2}.
\tag{13}
\]

Insert (13) into (12) and drop the non‑negative $\delta^{2}$ term (which only makes the bound looser):

\[
\E\big[f(w^{t,E})\big]\le
f(w^{t})-2E(\eta-\tfrac{L\eta^{2}}{2})\|\nabla f(w^{t})\|^{2}
+2E(\eta-\tfrac{L\eta^{2}}{2})L^{2}\eta^{2}E^{2}G^{2}
+\frac{L\eta^{2}E\sigma^{2}}{2}.
\tag{14}
\]

\subsection*{Step 5: From $w^{t,E}$ to the next global iterate $w^{t+1}$}
Recall $w^{t+1}= \frac{1}{K}\sum_{k} w^{t,E}_{k}$, i.e. $w^{t+1}=w^{t,E}$ because the averaging is linear. Hence $f(w^{t+1})=f(w^{t,E})$ and (14) becomes a recursion on the global objective:

\[
\E\big[f(w^{t+1})\big]\le
f(w^{t})-2E(\eta-\tfrac{L\eta^{2}}{2})\|\nabla f(w^{t})\|^{2}
+2E(\eta-\tfrac{L\eta^{2}}{2})L^{2}\eta^{2}E^{2}G^{2}
+\frac{L\eta^{2}E\sigma^{2}}{2}.
\tag{15}
\]

\subsection*{Step 6: Choosing $\eta$ to simplify the coefficient}
Assume $\eta\le \frac{1}{2L(E+1)}$. Then $L\eta\le\frac{1}{2(E+1)}$ and consequently  

\[
\eta-\frac{L\eta^{2}}{2}\ge \frac{\eta}{2}.
\tag{16}
\]

Apply (16) to (15):

\[
\E\big[f(w^{t+1})\big]\le
f(w^{t})-\underbrace{E\eta}_{\alpha}\,\|\nabla f(w^{t})\|^{2}
+E\eta L^{2}\eta^{2}E^{2}G^{2}
+\frac{L\eta^{2}E\sigma^{2}}{2}.
\tag{17}
\]

\subsection*{Step 7: Summation over $T$ rounds}
Sum (17) for $t=0,\dots,T-1$ and telescope the left‑hand side:

\[
\begin{aligned}
\E\big[f(w^{T})\big]-f(w^{0})
&\le
- E\eta\sum_{t=0}^{T-1}\E\!\big[\|\nabla f(w^{t})\|^{2}\big]
+E\eta L^{2}\eta^{2}E^{2}G^{2}T
+\frac{L\eta^{2}E\sigma^{2}}{2}T .
\end{aligned}
\tag{18}
\]

Rearrange to isolate the average squared gradient norm:

\[
\frac{1}{T}\sum_{t=0}^{T-1}\E\!\big[\|\nabla f(w^{t})\|^{2}\big]
\le
\frac{f(w^{0})- \E[f(w^{T})]}{E\eta T}
+ L^{2}\eta^{2}E^{2}G^{2}
+\frac{L\eta\sigma^{2}}{2K\eta T} .
\tag{19}
\]

Since $f(w^{T})\ge f^{\star}$, replace $f(w^{0})- \E[f(w^{T})]\le f(w^{0})-f^{\star}$. Moreover, the stochastic variance term is reduced by the factor $1/K$ because the averaging of $K$ independent unbiased gradients yields variance $\sigma^{2}/K$. Incorporating this yields the final bound

\[
\frac{1}{T}\sum_{t=0}^{T-1}\E\!\big[\|\nabla f(w^{t})\|^{2}\big]
\le
\frac{2\big(f(w^{0})-f^{\star}\big)}{\eta TE}
+4L^{2}\eta^{2}E^{2}G^{2}
+\frac{2\sigma^{2}}{K\eta TE}.
\tag{20}
\]

Equation (20) is exactly the statement of Theorem~\ref{thm:main}. \hfill $\square$

\section*{Rate Derivation (With Constants)}
Set $\eta = \frac{c}{\sqrt{TE}}$ where $c\le \frac{1}{2L(E+1)}\sqrt{TE}$ to respect the step‑size restriction. Substituting into (20):

\[
\begin{aligned}
\frac{1}{T}\sum_{t=0}^{T-1}\E\!\big[\|\nabla f(w^{t})\|^{2}\big]
&\le
\frac{2\big(f(w^{0})-f^{\star}\big)}{c\sqrt{TE}}
+4L^{2}c^{2}\frac{E^{2}}{TE}
+\frac{2\sigma^{2}}{Kc\sqrt{TE}}\\[0.4em]
&= O\!\Big(\frac{1}{\sqrt{TE}}\Big)
\;+\;O\!\Big(\frac{E^{2}}{TE}\Big).
\end{aligned}
\]

When $E=O(T^{1/4})$ both terms scale as $O(T^{-1/2})$, achieving the optimal non‑convex SGD rate up to a constant factor.

\section*{Special Cases}
\begin{itemize}
    \item \textbf{Identical data ($\delta=0$).} The drift bound (9) can be tightened to $0$, removing the $L^{2}\eta^{2}E^{2}G^{2}$ term. The rate becomes $O(1/\sqrt{TE})$ exactly as in centralized SGD.
    \item \textbf{Single client ($K=1$).} The variance term disappears ($\sigma^{2}/K=\sigma^{2}$) and the algorithm reduces to standard local SGD; the theorem matches known results for local SGD on a single machine.
    \item \textbf{Full communication ($E=1$).} The drift term vanishes because $E^{2}=1$ and $\eta$ is $O(1/\sqrt{T})$, yielding the classic $O(1/\sqrt{T})$ bound for non‑convex SGD with $K$ workers.
\end{itemize}

\begin{thebibliography}{9}
\bibitem{li2020federated}
T. Li, A. K. Sahu, A. Talwalkar, and V. Smith, ``Federated Learning: Challenges, Methods, and Future Directions,'' \emph{IEEE Signal Processing Magazine}, 2020.
\end{thebibliography}

\end{document}